% BleModule
\section{【応用】BleModule — スレッドセーフティパターン}

\begin{patternbox}{BleModule で学べるパターン}
\begin{itemize}
    \item volatileフラグ方式による安全なデータ受け渡し
    \item BLEコールバック（別タスク）→内部バッファ→update()→SystemData
    \item deinit()によるBLEリソースの解放
    \item 双方向通信（RX: コールバック受信、TX: update()で送信）
\end{itemize}
\end{patternbox}

\subsection{内部バッファとvolatileフラグ}

\begin{lstlisting}[style=cppstyle, caption={BleModule.h — 内部バッファ}]
class BleModule : public IModule {
private:
    // 内部バッファ（BLEコールバックスレッドが書き込む）
    uint8_t       _rxBuffer[BLE_RX_BUFFER_SIZE] = {};
    uint8_t       _rxBufferLength = 0;
    volatile bool _dataReceived   = false;  // volatileフラグ
    volatile bool _clientConnected = false;
public:
    // BLEコールバックから呼ばれる（別タスクで実行される）
    void onReceive(const uint8_t* payload, uint8_t length);
    // ...
};
\end{lstlisting}

\begin{cautionbox}[なぜvolatileが必要か]
BLEコールバックはArduinoのメインループとは別のFreeRTOSタスクで実行される。
\texttt{volatile}がないと、コンパイラが\texttt{\_dataReceived}の値をレジスタにキャッシュし、
メインループ側で変更を検出できない場合がある。
\texttt{volatile}は「この変数は外部から変更されるかもしれない」とコンパイラに伝えるキーワード。
\end{cautionbox}

\subsection{コールバック→内部バッファ→SystemData}

\begin{lstlisting}[style=cppstyle, caption={BleModule.cpp — データ受け渡しの流れ}]
// BLEコールバック（別タスクから呼ばれる）
void BleModule::onReceive(const uint8_t* payload, uint8_t length) {
    memcpy(_rxBuffer, payload, length);  // (1) バッファに書き込み
    _rxBufferLength = length;
    _dataReceived = true;                // (2) フラグを立てる
}

// update()は3フェーズループ内で呼ばれる（メインタスク）
void BleModule::update(SystemData& data) {
    data.ble.connected = _clientConnected;

    if (_dataReceived) {                             // (3) フラグ確認
        memcpy(data.ble.rxData, _rxBuffer, _rxBufferLength); // (4) コピー
        data.ble.rxLength  = _rxBufferLength;
        data.ble.newDataIn = true;
        _dataReceived = false;                       // (5) フラグクリア
    }

    // 送信リクエスト処理
    if (data.ble.sendRequest && _clientConnected && _txChar) {
        _txChar->setValue(data.ble.txData, data.ble.txLength);
        _txChar->notify();
        data.ble.sendRequest = false;
    }
}
\end{lstlisting}

\begin{pointbox}[volatileフラグ方式の安全性]
\texttt{bool}型の読み書きは原子的（1命令で完了する）。
「バッファに全データを書き込んだ後にフラグを立てる」→「フラグを確認してからバッファを読む」
の順序により、バッファの途中状態を読むことがない。
mutexやセマフォは不要で、割り込みコンテキストでも安全に使える。
\end{pointbox}

\subsection{BLEコールバックの登録}

\begin{lstlisting}[style=cppstyle, caption={BleModule.cpp — コールバッククラス}]
// BLEライブラリのコールバッククラスを継承
class BleRxCallbacks : public BLECharacteristicCallbacks {
private:
    BleModule* _module;
public:
    BleRxCallbacks(BleModule* module) : _module(module) {}
    void onWrite(BLECharacteristic* characteristic) override {
        std::string value = characteristic->getValue();
        uint8_t len = min((int)value.length(), BLE_RX_BUFFER_SIZE);
        _module->onReceive(
            reinterpret_cast<const uint8_t*>(value.data()), len);
    }
};
\end{lstlisting}

\begin{pointbox}[コールバックからモジュールへの橋渡し]
BLEライブラリはコールバッククラスの継承を要求するため、
内部クラスでラップして\texttt{BleModule::onReceive()}に橋渡しする。
コールバッククラスはモジュールのポインタを保持し、受信データをモジュールの内部バッファに渡す。
\end{pointbox}

\subsection{送信の流れ（ロジックフェーズ→update()）}

\begin{lstlisting}[style=cppstyle, caption={ロジックフェーズでの送信リクエスト}]
void applyPattern(SystemData& data) {
    // BLE受信データの処理
    if (data.ble.newDataIn) {
        // 受信データを処理...
        data.ble.newDataIn = false;  // 処理済みフラグクリア
    }

    // BLE送信（応答を返す例）
    if (needToRespond) {
        memcpy(data.ble.txData, response, responseLen);
        data.ble.txLength = responseLen;
        data.ble.sendRequest = true;  // 次のupdate()で送信される
    }
}
\end{lstlisting}

\begin{notebox}[双方向通信のデータフロー]
\textbf{受信}: BLEコールバック → 内部バッファ → update() → SystemData.ble.rxData → ロジックフェーズ\\
\textbf{送信}: ロジックフェーズ → SystemData.ble.txData → update() → BLE Notify\\
受信は入力フェーズの流れ、送信は出力フェーズの流れに対応する。
BleModuleは入力配列に登録するが、送信処理も\texttt{update()}内で行うため実質的に入出力兼用。
\end{notebox}
